\documentclass[10pt,a4paper]{amsart}
\usepackage{amsmath,amssymb,amsthm,graphicx}
\usepackage{verbatim}
\newenvironment{reference}{\comment}{\endcomment}
%\newenvironment{reference}{}{}
\newenvironment{slogan}{\comment}{\endcomment}
\newenvironment{history}{\comment}{\endcomment}

% For commutative diagrams we use Xy-pic
\usepackage[all]{xy}

% We use 2cell for 2-commutative diagrams.
\xyoption{2cell}
% Theorem environments.
%
\theoremstyle{plain}
\newtheorem{theorem}[subsection]{Theorem}
\newtheorem{proposition}[subsection]{Proposition}
\newtheorem{lemma}[subsection]{Lemma}

\theoremstyle{definition}
\newtheorem{definition}[subsection]{Definition}
\newtheorem{example}[subsection]{Example}
\newtheorem{exercise}[subsection]{Exercise}
\newtheorem{situation}[subsection]{Situation}

\theoremstyle{remark}
\newtheorem{remark}[subsection]{Remark}
\newtheorem{remarks}[subsection]{Remarks}

\numberwithin{equation}{subsection}

% Macros
%
\def\lim{\mathop{\mathrm{lim}}\nolimits}
\def\colim{\mathop{\mathrm{colim}}\nolimits}
\def\Spec{\mathop{\mathrm{Spec}}}
\def\Hom{\mathop{\mathrm{Hom}}\nolimits}
\def\Ext{\mathop{\mathrm{Ext}}\nolimits}
\def\SheafHom{\mathop{\mathcal{H}\!\mathit{om}}\nolimits}
\def\SheafExt{\mathop{\mathcal{E}\!\mathit{xt}}\nolimits}
\def\Sch{\mathit{Sch}}
\def\Mor{\mathop{\mathrm{Mor}}\nolimits}
\def\Ob{\mathop{\mathrm{Ob}}\nolimits}
\def\Sh{\mathop{\mathit{Sh}}\nolimits}
\def\NL{\mathop{N\!L}\nolimits}
\def\CH{\mathop{\mathrm{CH}}\nolimits}
\def\proetale{{pro\text{-}\acute{e}tale}}
\def\etale{{\acute{e}tale}}
\def\QCoh{\mathit{QCoh}}
\def\Ker{\mathop{\mathrm{Ker}}}
\def\Im{\mathop{\mathrm{Im}}}
\def\Coker{\mathop{\mathrm{Coker}}}
\def\Coim{\mathop{\mathrm{Coim}}}

% Boxtimes
%
\DeclareMathSymbol{\boxtimes}{\mathbin}{AMSa}{"02}

%
% Macros for moduli stacks/spaces
%
\def\QCohstack{\mathcal{QC}\!\mathit{oh}}
\def\Cohstack{\mathcal{C}\!\mathit{oh}}
\def\Spacesstack{\mathcal{S}\!\mathit{paces}}
\def\Quotfunctor{\mathrm{Quot}}
\def\Hilbfunctor{\mathrm{Hilb}}
\def\Curvesstack{\mathcal{C}\!\mathit{urves}}
\def\Polarizedstack{\mathcal{P}\!\mathit{olarized}}
\def\Complexesstack{\mathcal{C}\!\mathit{omplexes}}
% \Pic is the operator that assigns to X its picard group, usage \Pic(X)
% \Picardstack_{X/B} denotes the Picard stack of X over B
% \Picardfunctor_{X/B} denotes the Picard functor of X over B
\def\Pic{\mathop{\mathrm{Pic}}\nolimits}
\def\Picardstack{\mathcal{P}\!\mathit{ic}}
\def\Picardfunctor{\mathrm{Pic}}
\def\Deformationcategory{\mathcal{D}\!\mathit{ef}}
\title{A group scheme whose identity component is flat but which is not flat}
\author{The Stacks Project Authors}
\date{stacks-project:examples.tex:local-snapshot-20261012; GNU FDL 1.2 or later}
\begin{document}
\maketitle
\noindent\textit{Editorial note.} All delivered mathematical text is the author's own, copied from the pinned Stacks Project chapter file examples.tex; no statement or proof step was rewritten, reordered or removed.\par\medskip
\noindent\textit{Editorial note.} The counted proof is the author's complete argument: The monomorphism identities X x\_S X = X and X x\_S X x\_S X = X with unique S-morphisms are the author's own carried step (3905-3908); the identity component of S amal X is the S summand.\par\medskip
\noindent\textit{Editorial note.} The label lemma-non-flat-group-scheme names the counted statement, and the proof environment immediately following it in the source is delivered verbatim as its proof.\par\medskip
\noindent\textit{Editorial note.} The counted statement is carried by the author's own lemma environment in the chapter source.\par\medskip
\noindent\textit{Allowed standard prerequisites (editorial).}\begin{itemize}
\item Context span 3890-3897: for a monomorphism of schemes X -> S, set G = S amal X and define m: G x\_S G -> G on the decomposition G x\_S G = (X x\_S X) amal X amal X amal S by mapping the summands X x\_S X and S into S and the summands X into X by the identity
\item Context span 3898-3909: the group-law associativity m compose (m x id\_G) = m compose (id\_G x m) is verified by decomposing G x\_S G x\_S G into the 8 pieces, each isomorphic to S, X, X x\_S X, or X x\_S X x\_S X; since X -> S is a monomorphism, X x\_S X = X and X x\_S X x\_S X = X and in each case there is exactly one S-morphism S -> S or X -> X, whence the two compositions agree
\item Context span 3910-3914: taking X -> S any nonflat monomorphism (e.g., a closed immersion) gives a group scheme over S whose identity component is S (hence flat) but which is not flat
\end{itemize}\medskip
Let $X \to S$ be a monomorphism of schemes. Let $G = S \amalg X$.
Let $m : G \times_S G \to G$ be the $S$-morphism
$$
G \times_S G = X \times_S X \amalg X \amalg X \amalg S
\longrightarrow G = X \amalg S
$$
which maps the summands $X \times_S X$ and $S$ into $S$ and
maps the summands $X$ into $X$ by the identity morphism.
This defines a group law. To see this we have to show that
$m \circ (m \times \text{id}_G) = m \circ (\text{id}_G \times m)$
as maps $G \times_S G \times_S G \to G$. Decomposing
$G \times_S G \times_S G$ into components as above, we see that
we need to verify this for the restriction to each of the $8$-pieces.
Each piece is isomorphic to either $S$, $X$, $X \times_S X$, or
$X \times_S X \times_S X$. Moreover, both maps map these pieces
to $S$, $X$, $S$, $X$ respectively. Having said this, the fact that
$X \to S$ is a monomorphism implies that $X \times_S X \cong X$
and $X \times_S X \times_S X \cong X$ and that there is in each case
exactly one $S$-morphism $S \to S$ or $X \to X$. Thus we see
that $m \circ (m \times \text{id}_G) = m \circ (\text{id}_G \times m)$.
Thus taking $X \to S$ to be any
nonflat monomorphism of schemes (e.g., a closed immersion)
we get an example of a group scheme
over a base $S$ whose identity component is $S$ (hence flat)
but which is not flat.

\begin{lemma}
\label{lemma-non-flat-group-scheme}
There exists a group scheme $G$ over a base $S$ whose identity
component is flat over $S$ but which is not flat over $S$.
\end{lemma}

\begin{proof}
See discussion above.
\end{proof}

\end{document}
